\subsection{FRM Exam 2003: VaR and Expected Shortfall (90\%) from ordered returns}

\textbf{Enunciado}

\emph{Given the following 30 ordered percentage returns of an asset, calculate the VAR and expected shortfall at a 90\% confidence level:}
\[
-16,-14,-10,-7,-7,-5,-4,-4,-4,-3,-1,-1,0
\]
\[
0,0,1,2,2,4,6,7,8,9,11,12,12,14,18,21,23.
\]
\begin{enumerate}[label=\alph*.]
\item VaR (90\%) = 10, expected shortfall = 14
\item \hl{VaR (90\%) = 10, expected shortfall = 15}
\item VaR (90\%) = 14, expected shortfall = 15
\item VaR (90\%) = 18, expected shortfall = 22
\end{enumerate}

\subsubsection*{Solución}


\textbf{VaR(90\%) por simulación histórica.} Los retornos están ordenados de peor a mejor. Un nivel de confianza 90\% equivale a una cola izquierda del 10\%.
Con \(n=30\):
\[
0.10\times 30=3.
\]
El 3er peor retorno es \(-10\%\). Entonces el VaR(90\%) como magnitud de pérdida es:
\[
\VaR_{90\%}\approx 10.
\]

\textbf{ES(90\%) como promedio en la cola.} Bajo la convención habitual de examen ``más allá del VaR'' (estrictamente peores que \(-10\%\)):
\[
\{-16\%,-14\%\}.
\]
El ES(90\%) como magnitud de pérdida es:
\[
\ES_{90\%}\approx \frac{16+14}{2}=15.
\]


Con esto, \(\VaR=10\) y \(\ES=15\), así que la opción (b) es correcta.












